\chapter{Observateur couplé, auto-inertie et conservation généalogique de l'information}
\label{chap:observador-autoinercia-landauer}

\section{État conjoint et résidence de l'observateur}

La description dynamique élémentaire comprend la section observée,
l'observateur réalisé et le registre qu'ils produisent tous deux. Nous écrivons
\begin{equation}
 \boxed{\Gamma=(x,\mathcal O,m)\in
 \mathcal X_{\rm TPK}\times\mathfrak O\times\mathfrak M.}
 \label{eq:rm-obs-joint-state}
\end{equation}
La coordonnée \(x\) conserve cylindre, chemin, feuillet, orientation, retenue,
frontière et possibilités de prolongement ; \(\mathcal O\) déclare le
sous-système lecteur et sa résolution ; \(m\) conserve la mémoire distinguable
du couplage. L'état conjoint constitue le domaine primaire de la
mesure et de la réponse inertielle.

Un observateur à la profondeur \(n\) est représenté au moyen d'une carte
\begin{equation}
 q_{\mathcal O,n}:\mathcal X_{\rm TPK}^{(n)}
 \longrightarrow Y_{\mathcal O,n}.
 \label{eq:rm-obs-observer-chart}
\end{equation}
Pour \(x\in\mathcal X_{\rm TPK}^{(n)}\), son horizon observationnel est la
fibre
\begin{equation}
 \mathfrak H_{\mathcal O,n}(x)
 =q_{\mathcal O,n}^{-1}(q_{\mathcal O,n}(x)).
 \label{eq:rm-obs-horizon-fiber}
\end{equation}
La fibre réunit les sections indiscernables pour la carte déclarée. Le
raffinement conserve la naturalité lorsqu'il existe des applications
\(r_{n+1,n}\) et \(s_{n+1,n}\) telles que
\begin{equation}
 q_{\mathcal O,n}\circ r_{n+1,n}
 =s_{n+1,n}\circ q_{\mathcal O,n+1}.
 \label{eq:rm-obs-observer-naturality}
\end{equation}
Cette identité rend l'observation finie compatible avec la généalogie
profonde : le résultat de la troncature d'une lecture raffinée coïncide avec la
lecture de l'état tronqué.

\begin{definition}[Observateur réalisé]
Un observateur réalisé est un tuple
\begin{equation}
 \mathcal O_a=(m_0,U_a,q_a,\mathcal C_a)
 \label{eq:rm-obs-realized-observer}
\end{equation}
formé d'une préparation \(m_0\), d'un couplage \(U_a\), d'un lecteur
de pointeur \(q_a\) et d'une condition de compatibilité et de prolongement
\(\mathcal C_a\). Sa fonction consiste à produire un registre distinguable
et à l'employer comme condition de continuation du chemin.
\end{definition}

La définition situe le cadre de lecture à l'intérieur du système. La
coordonnée publiée dépend du couplage, de la résolution et de la
mémoire ; la section enrichie conserve la totalité dont provient cette
coordonnée.

\section{Contexte de mesure et conditionnement par fibres}

Soient \((\Omega_S,\Sigma_S)\) et \((\Omega_M,\Sigma_M)\) des espaces mesurables
de sections du système et de l'appareil. Le contexte \(a\) emploie un
couplage mesurable
\begin{equation}
 U_a:\Omega_S\times\Omega_M
 \longrightarrow\Omega_S\times\Omega_M
 \label{eq:rm-obs-measurable-coupling}
\end{equation}
et un lecteur \(q_a:\Omega_M\to Y_a\). La sortie est
\begin{equation}
 \operatorname{Out}_a(x)
 =q_a\!\left(\operatorname{pr}_M U_a(x,m_0)\right).
 \label{eq:rm-obs-output}
\end{equation}
Pour un résultat \(y\), on définit la fibre mesurable
\begin{equation}
 \Omega_{a,y}
 =\{x\in\Omega_a:\operatorname{Out}_a(x)=y\}.
 \label{eq:rm-obs-outcome-fiber}
\end{equation}
La mise à jour généalogique est la restriction
\begin{equation}
 \Omega_a\longmapsto\Omega_{a,y}.
 \label{eq:rm-obs-genealogical-update}
\end{equation}
À profondeur finie, l'opération conserve les préfixes qui admettent une
extension dans \(\Omega_{a,y}\). La mémoire du résultat s'incorpore à la
condition de prolongement ; le chemin parcouru reste enregistré.

Si \(\mu_a(\Omega_{a,y})>0\), la mesure conditionnée est
\begin{equation}
 \mu_{a,y}(B)
 =\frac{\mu_a(B\cap\Omega_{a,y})}
 {\mu_a(\Omega_{a,y})}.
 \label{eq:rm-obs-conditional-measure}
\end{equation}
Cette expression constitue la carte probabiliste de l'état relatif.
Sa réalisation spectrale s'obtient au moyen de projections cylindriques.

\section{PVM cylindrique et correspondance de Lüders}

Supposons que \(\Omega_{a,y}\) est une union finie de cylindres
disjoints d'une profondeur commune \(n\). Soit \(P_y\) la somme de leurs
projecteurs orthogonaux dans la réalisation de Hilbert. Pour un état
de densité \(\rho\) et une PVM commune d'événements cylindriques, la probabilité
de Born est
\begin{equation}
 p(y)=\Tr(\rho P_y).
 \label{eq:rm-obs-born}
\end{equation}
Lorsque \(p(y)>0\), la mise à jour de Lüders \cite{Luders1951} est
\begin{equation}
 \rho_y=\frac{P_y\rho P_y}{\Tr(\rho P_y)}.
 \label{eq:rm-obs-luders}
\end{equation}

\begin{theorem}[Équivalence entre restriction de fibre et Lüders]
\label{thm:rm-obs-fiber-luders}
Soit \(\mathcal A_n\) l'algèbre d'événements cylindriques à la profondeur
\(n\). Supposons que
\begin{equation}
 \mu_a(B)=\Tr(\rho P_B),
 \qquad B\in\mathcal A_n,
 \label{eq:rm-obs-measure-pvm}
\end{equation}
et que \(P_B\) et \(P_y\) appartiennent à la même mesure spectrale projective.
Alors
\begin{equation}
 \boxed{
 \Tr(\rho_yP_B)
 =\frac{\mu_a(B\cap\Omega_{a,y})}
 {\mu_a(\Omega_{a,y})}.}
 \label{eq:rm-obs-luders-fiber-equivalence}
\end{equation}
\end{theorem}

\begin{proof}
L'appartenance à une PVM commune donne
\(P_BP_y=P_{B\cap\Omega_{a,y}}\). La cyclicité de la trace et
\cref{eq:rm-obs-measure-pvm} produisent
\[
 \Tr(\rho_yP_B)
 =\frac{\Tr(\rho P_yP_BP_y)}{\Tr(\rho P_y)}
 =\frac{\Tr(\rho P_{B\cap\Omega_{a,y}})}
 {\Tr(\rho P_y)},
\]
ce qui coïncide avec \cref{eq:rm-obs-conditional-measure}.
\end{proof}

L'hypothèse de PVM commune est substantielle. Deux contextes mutuellement
non biaisés admettent leurs propres descriptions projectives et nécessitent une POVM
bruitée ou une dilatation pour une lecture conjointe. La correspondance de
\cref{thm:rm-obs-fiber-luders} est affirmée à l'intérieur du contexte spectral
déclaré.

\section{Instrument quantique, dilatation et conservation du canal}

Soit \(\widehat U_a\) la réalisation unitaire du couplage entre
système et appareil, et soit \(|m_0\rangle\) la préparation du pointeur. Une
base raffinée \(|m_{y,\alpha}\rangle\) induit les opérateurs de Kraus
\begin{equation}
 M_{y,\alpha}
 =\langle m_{y,\alpha}|\widehat U_a|m_0\rangle,
 \qquad
 \sum_{y,\alpha}M_{y,\alpha}^{\dagger}M_{y,\alpha}=I.
 \label{eq:rm-obs-kraus}
\end{equation}
L'instrument et son effet sont
\begin{equation}
 \mathcal I_y(\rho)
 =\sum_\alpha M_{y,\alpha}\rho M_{y,\alpha}^{\dagger},
 \qquad
 E_y=\sum_\alpha M_{y,\alpha}^{\dagger}M_{y,\alpha}.
 \label{eq:rm-obs-instrument}
\end{equation}
La probabilité et l'état relatif prennent la forme
\begin{equation}
 p(y)=\Tr(\rho E_y),
 \qquad
 \rho_y=\frac{\mathcal I_y(\rho)}{p(y)}.
 \label{eq:rm-obs-general-conditioned-state}
\end{equation}
Le canal non sélectif est
\begin{equation}
 \mathcal E_a(\rho)=\sum_y\mathcal I_y(\rho).
 \label{eq:rm-obs-unconditioned-channel}
\end{equation}

\begin{proposition}[Conservation exacte du canal global]
\label{prop:rm-obs-global-channel}
L'isométrie
\begin{equation}
 V_a:\mathcal H_S\longrightarrow\mathcal H_S\otimes\mathcal H_M,
 \qquad
 V_a|\psi\rangle=\widehat U_a(|\psi\rangle\otimes|m_0\rangle)
 \label{eq:rm-obs-stinespring-isometry}
\end{equation}
satisfait \(V_a^\dagger V_a=I\). En conséquence,
\begin{equation}
 \sum_y p(y)=1,
 \qquad
 \Tr\mathcal E_a(\rho)=\Tr\rho,
 \label{eq:rm-obs-trace-preservation}
\end{equation}
et \(\mathcal E_a\) est complètement positif et préserve la trace.
Pour un état pur initial, l'état conjoint
\(V_a|\psi\rangle\) reste pur et normalisé.
\end{proposition}

\begin{proof}
L'unitarité de \(\widehat U_a\) et la normalisation de
\(|m_0\rangle\) donnent \(V_a^\dagger V_a=I\). L'identité de complétude de
\cref{eq:rm-obs-kraus} préserve la trace. La forme de Kraus démontre
la positivité complète, et une isométrie envoie les vecteurs unitaires sur des
vecteurs unitaires.
\end{proof}

La conservation du canal global signifie que l'évolution dilatée
retient toutes les amplitudes et tous les registres. Le conditionnement
sélectionne une fibre relative pour le lecteur qui a obtenu \(y\). La somme
des instruments reproduit le canal complet. Cette triple distinction
sépare évolution, lecture et état relatif.

Le théorème de Stinespring \cite{Stinespring1955} garantit une dilatation isométrique pour tout
canal complètement positif. L'identification de l'espace auxiliaire à
la mémoire TPK se réalise au moyen d'un entrelaceur qui conserve chemin,
retenue, registre et raffinement. L'existence abstraite de la dilatation
et la généalogie concrète du registre sont deux résultats coordonnables et
typés.

\section{État relatif et formulation everettienne}

La formulation de l'état relatif fournit l'antécédent conceptuel
de cette construction \cite{Everett1957}.

Dans une prémesure idéale, supposons que les états de pointeur
\(\{|m_r\rangle\}_r\) sont orthonormaux et que \(\widehat U_a\) est une
isométrie sur le sous-espace engendré par
\(\{|r\rangle\otimes|m_0\rangle\}_r\). Alors
\begin{equation}
 \widehat U_a(|r\rangle\otimes|m_0\rangle)
 =|r\rangle\otimes|m_r\rangle.
 \label{eq:rm-obs-ideal-premeasurement}
\end{equation}
Pour \(|\psi\rangle=\sum_r c_r|r\rangle\), la linéarité produit
\begin{equation}
 |\Psi'\rangle
 =\sum_r c_r|r\rangle\otimes|m_r\rangle.
 \label{eq:rm-obs-entangled-record}
\end{equation}
Le registre \(m_r\) définit la section relative
\begin{equation}
 \rho_r
 =\frac{Q_r|\Psi'\rangle\langle\Psi'|Q_r}
 {\Tr(|\Psi'\rangle\langle\Psi'|Q_r)},
 \qquad
 Q_r=I_S\otimes|m_r\rangle\langle m_r|.
 \label{eq:rm-obs-relative-state}
\end{equation}
Dans l'espace généalogique, le même résultat correspond à
\begin{equation}
 \Omega_{a,r}
 =\{x:\operatorname{Out}_a(x)=r\}.
 \label{eq:rm-obs-relative-fiber}
\end{equation}

\begin{theorem}[Correspondance d'état relatif]
\label{thm:rm-obs-everett-correspondence}
Sous les hypothèses projectives de
\cref{thm:rm-obs-fiber-luders}, l'état relatif de
\cref{eq:rm-obs-relative-state} et la mesure conditionnée sur
\(\Omega_{a,r}\) attribuent les mêmes probabilités à tous les événements de
la PVM cylindrique déclarée. L'isométrie globale conserve la superposition
corrélée de tous les registres.
\end{theorem}

\begin{proof}
La première affirmation est
\cref{eq:rm-obs-luders-fiber-equivalence}. La deuxième découle de
\(V_a^\dagger V_a=I\) et de l'expression linéaire
\cref{eq:rm-obs-entangled-record}.
\end{proof}

La correspondance everettienne atteint le niveau mathématique des états
relatifs : l'état conjoint conserve les canaux et chaque registre détermine
une fibre conditionnée. HMT ajoute à chaque fibre une généalogie composée
de préfixes, de retenue, d'orientation, de frontière et de prolongements compatibles.
La stabilité macroscopique des registres appartient à la théorie de
la décohérence de l'appareil ; l'ontologie de branches physiquement autonomes
appartient à une interprétation supplémentaire. Le théorème précédent conserve sa
force indépendamment de ce choix ontologique.

